\documentclass[10pt]{article}

\usepackage[margin=0.75in]{geometry}
\usepackage{amsmath,amsthm,amssymb,bm}
\usepackage{xcolor}
\usepackage{cancel}
\usepackage{graphicx}
\usepackage{changepage}
\usepackage{circuitikz}
\usepackage{pgfplots}
\usepackage{minted}
\usepackage{physics}
\usepackage{tikz}
\usepackage{tikz-3dplot}
\usepackage{hyperref}
\usepackage{siunitx}
\usepackage{fontspec}
\usepackage{relsize}
\usepackage{subfig}
\usepackage{todonotes}
\usepackage{contour}
\usepackage{multicol, multirow, booktabs}
\usepackage[breakable]{tcolorbox}
\usepackage[inline]{enumitem}

\theoremstyle{definition}
\newtheorem{problem}{Problem}
\newtheorem{soln}{Solution}

\pgfplotsset{compat=newest}
\usetikzlibrary{lindenmayersystems}
\usetikzlibrary{arrows}
\usetikzlibrary{calc}
\usetikzlibrary{positioning, fit}
\usetikzlibrary{3d, perspective}
\usetikzlibrary{math}
\usetikzlibrary{fadings}
\usetikzlibrary{angles,quotes} 
\usetikzlibrary{decorations.markings}
\usetikzlibrary{tikzmark}

\definecolor{incolor}{HTML}{303F9F}
\definecolor{outcolor}{HTML}{D84315}
\definecolor{cellborder}{HTML}{CFCFCF}
\definecolor{cellbackground}{HTML}{F7F7F7}
\newcommand{\ui}{\hat{i}}
\newcommand{\uj}{\hat{j}}
\newcommand{\uk}{\hat{k}}
\newcommand{\ux}{\hat{\mathbf{x}}}
\newcommand{\uy}{\hat{\mathbf{y}}}
\newcommand{\uz}{\hat{\mathbf{z}}}
\newcommand{\uv}[1]{\hat{\bv{{#1}}}}
\newcommand{\pr}[1]{{#1^\prime}}
\newcommand{\pd}[2][]{{\frac{\partial^{#1}}{\partial {{#2}^{#1}}}}}
\newcommand{\dif}[2][]{{\frac{d^{#1}}{d{{#2}^{#1}}}}}
\newcommand{\grd}{\nabla}

\pgfdeclarelayer{background}  
\pgfsetlayers{background,main}
\AtBeginDocument{\RenewCommandCopy\qty\SI}

\makeatletter
\newcommand{\boxspacing}{\kern\kvtcb@left@rule\kern\kvtcb@boxsep}
\makeatother
\newcommand{\prompt}[4]{
    \ttfamily\llap{{\color{#2}[#3]:\hspace{3pt}#4}}\vspace{-\baselineskip}
}

\newcommand{\thevenin}[2]{
  \begin{center}
    \begin{circuitikz} \draw
      (0,0) -- (2,0) to[battery1, l_=$V_{Th}\eq#1$] (2,2) 
      to[resistor, l_=$R_{Th}\eq#2$] (0,2)
      ;
      \draw [o-] (-.07,2.079);
      \draw [o-] (-.07,0.079);
    \end{circuitikz}
  \end{center}
}

\newcommand{\norton}[2]{
  \begin{center}
    \begin{circuitikz} \draw
      (0,0) -- (3,0) to[american current source, l_=$I_{N}\eq#1$] (3,2) -- (0,2) (2,0)
      to[resistor, l=$R_{N}\eq#2$] (2,2)
      ;
      \draw [o-] (-.07,2.079);
      \draw [o-] (-.07,0.079);
    \end{circuitikz}
  \end{center}
}

\DeclareMathOperator{\Div}{div}
\DeclareMathOperator{\Curl}{curl}
\DeclareMathOperator{\sgn}{sgn}
\DeclareMathOperator{\erfc}{erfc}

\newcommand{\highlight}[1]{\colorbox{yellow}{$\displaystyle #1$}}
\newcommand{\ti}[1]{\widetilde{#1}}
\renewcommand{\div}[1]{\bv{\nabla}\cdot #1}
\renewcommand{\curl}[1]{\bv{\nabla}\times #1}

\newfontface{\Kaufmann}{Kaufmann}
\DeclareTextFontCommand{\kf}{\Kaufmann}
\newcommand{\scriptr}{\fontsize{12pt}{12pt}\kf{r}\fontsize{10pt}{10pt}}

\newfontface{\KaufmannB}{Kaufmann Bold}
\DeclareTextFontCommand{\kfb}{\KaufmannB}
\newcommand{\bscriptr}{\fontsize{12pt}{12pt}\kfb{r}\fontsize{10pt}{10pt}}
\newcommand{\justif}[2]{&{#1}&\text{#2}}
\newcommand{\bv}[1]{\bm{\mathrm{#1}}}

\title{Physics 4700H: Assignment IV}
\author{Jeremy Favro (0805980) \\ Trent University, Peterborough, ON, Canada}
\date{\today}

\begin{document}
\maketitle

% PROBLEM 1
\begin{problem}
Consider two systems of distinguishable noninteracting spin-$1/2$ particles in a uniform
magnetic field $\bv{B}$. Each spin has a magnetic moment along the field direction $\mu=\gamma\hbar m$ where
$m = \pm1/2$. The energy, $U$, of a system is proportional to its spin excess $X\equiv N_\uparrow-N_\downarrow$.
$$U=-\sum_{i=1}^{N}\bv{\mu}_i\cdot\bv{B}=-(N_\uparrow-N_\downarrow)\frac{\gamma\hbar}{2}B=-X\frac{\gamma\hbar}{2}B.$$
System 1: $N_1=8$ $X_{1i}=4$\\
System 2: $N_2=4$ $X_{2i}=-2$
Before the systems are put in thermal contact
\begin{enumerate}[label=(\alph*)]
  \item What are the initial multiplicity of System 1 ($\Omega_{1i}$), the initial multiplicity of System 2
        ($\Omega_{2i}$), and the initial multiplicity of the two systems together ($\Omega_{i}$)?
  \item What is an accessible microstate of this combined system before thermal contact?
\end{enumerate}
After the systems are put in thermal contact
\begin{enumerate}[label=(\alph*)]
  \setcounter{enumi}{2}
  \item What are $N$, $N_\uparrow$, and $N_\downarrow$ for the combined system?
  \item What is the spin excess, $X$, of the combined system?
  \item What is a microstate of the combined system which is only accessible after thermal
        contact?
  \item Calculate the multiplicity of the combined system, $\Omega_f$, using the binomial expansion.
  \item Calculate the multiplicity of the combined system, $\Omega_f$, using
        $$\Omega(N,X)=\sum_{X_1}\Omega_1(N_1,X_1)\Omega_2(N_2, X-X_1).$$
  \item What is the most probable value of $X_1$ in the combined system ($X_{1\text{max}}$)?
\end{enumerate}
\end{problem}
\begin{soln}~
  \begin{enumerate}[label=(\alph*)]
    \item For a binary system such as this one
          $$\Omega(N,X)=\frac{N!}{\left(\frac{1}{2}(N+X)\right)!\left(\frac{1}{2}(N-X)\right)!}$$
          and so $\Omega_{1i}=28$ and $\Omega_{2i}=4$ and $\Omega_{i}=\Omega_{1i}\Omega_{2i}=112$.
    \item One which preserves the associated $X$ values for systems 1 and 2. E.g. system 1 in
          $\uparrow\uparrow\uparrow\uparrow\uparrow\uparrow\downarrow\downarrow$ and system 2 in
          $\uparrow\downarrow\downarrow\downarrow$ and so the whole system is in
          $$\uparrow\uparrow\uparrow\uparrow\uparrow\uparrow\downarrow\downarrow\cdot \uparrow\downarrow\downarrow\downarrow.$$
    \item $N=N_1+N_2=12$, $N_\uparrow$ and $N_\downarrow$ can be determined by noting that, despite the above microstate having a
          specific ordering for the spins, any accessible microstate of the same energy would have to have the same number of up and down spins, hence
          $N_\uparrow=7$ and $N_\downarrow=5$.
    \item $X=N_\uparrow-N_\downarrow=7-5=2$.
    \item One which preserves the net $X=2$ but does not preserve the $X$ values of the individual systems, e.g.
          $$\uparrow\uparrow\uparrow\uparrow\uparrow\downarrow\downarrow\downarrow\cdot \uparrow\uparrow\downarrow\downarrow.$$
    \item We want to calculate how many ways we can have a system with $X=2$ for $N=12$. This is just choosing $N_\uparrow=7$
          spins to be up, and letting the rest be down, i.e.
          $$\begin{pmatrix}
              12 \\
              7
            \end{pmatrix}=792.$$
    \item As in class I will use a table to determine the allowed values $X_1$, i.e. those which give a possible value of $X_2=X-X_1$
          \begin{center}
            \begin{tabular}{c|c|c|c|c|c}
              $X_1$  & $X_2=X-X_1$ & $\Omega_1$ & $\Omega_2$ & $\Omega_f=\Omega_1\Omega_2$ & Legal? \\
              \cline{1-5}
              $8$    & $2-8=-6$    &            &            &                             & No     \\
              $6$    & $-4$        & 8          & 1          & 8                           & Yes    \\
              $4$    & $-2$        & 28         & 4          & 112                         & Yes    \\
              $2$    & $0$         & 56         & 6          & 336                         & Yes    \\
              $0$    & $2$         & 70         & 4          & 280                         & Yes    \\
              $-2$   & $4$         & 56         & 1          & 56                          & Yes    \\
              \vdots &             &            &            &                             & No
            \end{tabular}
          \end{center}
          and since anything with $X_1<-2$ will continue to increase $X_2$ past $4$, its highest possible value, no other $X_1$ values are valid and we end the table there.
          So the ``net'' multiplicity is the sum of the products, $8+112+336+280+56=792$, as we got in the previous part.
    \item The one with the highest multiplicity, given that all states are equally likely, $X_1=2$. Like in class, this is not
          the most likely state for the $X_1$ system in isolation, but when combined with the $X_2$ system, sufficient combined states are available
          as to increase its multiplicity past the $X_1=0$ state.
  \end{enumerate}
\end{soln}

% PROBLEM 2
\begin{problem}
In class, we considered two systems with $N_1$ and $N_2$ particles respectively where $N_1$ and
$N_2$ are large. Initially, they have energies $U_{1i}$ and $U_{2i}$. After they are brought into thermal
contact, $U = U_{1i} + U_{2i}$, and we found that $\Omega_1\Omega_2$ is maximum when $U_1 = U_{1\text{max}} = N_1U/N$.
In addition we considered the multiplicity when $U_1 = U_{1\text{max}} + \delta$, and we found that
$$\frac{\Omega(N,U,U_{1\text{max}}+\delta)}{\Omega(N,U,U_{1\text{max}})}=e^{-\alpha\delta^2/N_1}e^{-\alpha\delta^2/N_2}\text{ where }\alpha\equiv \frac{2}{(\gamma\hbar B)^2}.$$
This is not actually a probability but instead just a ratio of two values $\Omega_1\Omega_2$ corresponding
to two values of $U_1$. In this problem, let's sort out more precisely the probability of finding
the system in a state with $U_1 > U_{1\text{max}} \pm δ$ (that is in a state outside the
central peak). As in class, we'll use $N_1 = N_2 = 10^{22}$ and $\delta=10^{12}/\sqrt{\alpha}$.
The probability we want is the ratio of two values:
The numerator ($NUM$) is the sum of the multiplicities for all values of $U_1$ for which $U_1$ is
either greater than $U_{1\text{max}}+\delta$ or less than $U_{1\text{max}}-\delta$.
The denominator ($DEN$) is the sum of the multiplicities over all possible values of $U_1$.
\begin{enumerate}[label=(\alph*)]
  \item Find a (symbolic) expression for the denominator starting from
        $$DEN\approx \int_{-\infty}^{\infty}\Omega(N,U,U_{1\text{max}})e^{-\alpha u^2/N_1}e^{-\alpha u^2/N_2}\,du.$$
  \item Find a (symbolic) expression for the numerator starting from
        $$NUM\approx \int_{\delta}^{\infty}\Omega(N,U,U_{1\text{max}})e^{-\alpha u^2/N_1}e^{-\alpha u^2/N_2}\,du.$$
        Find a (symbolic) expression for this in terms of the complementary error function:
        $$\erfc(z)=\frac{2}{\sqrt{\pi}}\int_{z}^{\infty}e^{-t^2}dt.$$
  \item Now take the ratio $NUM/DEN$ and evaluate using the values of $N_1$, $N_2$, and $\delta$ given
        above.
\end{enumerate}
\end{problem}
\begin{soln}~
  \begin{enumerate}[label=(\alph*)]
    \item This is not a hard integral to do. First we combine the two exponentials (call the integral $I_1$)
          $$I_1=\Omega(N,U,U_{1\text{max}})\int_{-\infty}^{\infty}e^{-\alpha u^2\left[\frac{N_1+N_2}{N_1N_2}\right]}du$$
          and make the substitution $t=\sqrt{\alpha \beta} u\implies dt=\sqrt{\alpha\beta}du$ for $\beta=(N_1+N_2)/N_1N_2=\num{2e-22}$
          which transforms the integral to
          $$I_1=\frac{\Omega(N,U,U_{1\text{max}})}{\sqrt{\alpha\beta}}\int_{-\infty}^{\infty}e^{-t^2}dt$$
          where we need not change the bounds as $u=\pm\infty\implies t=\pm \infty$. This is just a standard Gaussian integral and so comes out to
          $\sqrt{\pi}$, yielding
          $$I_1=\frac{\Omega(N,U,U_{1\text{max}})}{\sqrt{\pi\alpha\beta}}$$
    \item Call the integral $I_2$. We again combine the exponentials and use the same definition for $\beta$
          $$I_2= 2\Omega(N,U,U_{1\text{max}})\int_{\delta}^{\infty}e^{-\alpha\beta u^2}du$$
          and we make the same substitution for $t=\sqrt{\alpha\beta}u$ transforming the integral to
          $$I_2=\frac{2\Omega(N,U,U_{1\text{max}})}{\sqrt{\alpha\beta}}\int_{\sqrt{\alpha\beta}\delta}^{\infty}e^{-t^2}dt$$
          where here we did need to change the lower bound as $\delta$ is finite. Using the given error function identity this integral ``evaluates'' to
          $$I_2=\frac{2\sqrt{\pi}\Omega(N,U,U_{1\text{max}})}{\sqrt{\alpha\beta}}\erfc(\sqrt{\alpha\beta}\delta).$$
    \item
          \begin{align*}
            \frac{NUM}{DEN} & \approx\frac{\frac{2\sqrt{\pi}\Omega(N,U,U_{1\text{max}})}{\sqrt{\alpha\beta}}\erfc(\sqrt{\alpha\beta}\delta)}{\frac{\Omega(N,U,U_{1\text{max}})}{\sqrt{\pi\alpha\beta}}}                    \\
                            & =\frac{2\sqrt{\pi}\cancel{\Omega(N,U,U_{1\text{max}})}\erfc(\sqrt{\alpha\beta}\delta)\sqrt{\pi}\cancel{\sqrt{\alpha\beta}}}{\cancel{\sqrt{\alpha\beta}}\cancel{\Omega(N,U,U_{1\text{max}})}} \\
                            & =2\pi\erfc(\sqrt{\alpha\beta}\delta)                                                                                                                                                         \\
                            & =2\pi(1-\erf(\sqrt{2}10^{-11}10^{12}))=2\pi(1-\erf(10\sqrt{2})).
          \end{align*}
          Plugging that into SageMath I obtain
          \begin{tcolorbox}[breakable, size=fbox, boxrule=1pt, pad at break*=1mm,colback=cellbackground, colframe=cellborder]
            \prompt{In}{incolor}{1}{\boxspacing}
            \begin{minted}[breaklines, autogobble]{sage}
      show(N(2*pi*(1-erf(sqrt(2)*10)), prec=512))
              \end{minted}
          \end{tcolorbox}
          \begin{tcolorbox}[breakable, size=fbox, boxrule=.5pt, pad at break*=1mm, opacityfill=0]
            \prompt{Out}{outcolor}{1}{\boxspacing}
            \begin{align*}
               & 3.460306120704405305148892561702650171704972399133244869217155326584265 \\
               & 25521418956578529376665302960845474399137567655796603836973700352396318 \\
               & 381498018335 \times 10^{-88}
            \end{align*}
          \end{tcolorbox}
          The \verb|N| function in Sage gives a numerical approximation to \verb|prec| bits of precision. Ordinary floating point precision, 53 bits for 
          the actual number, isn't enough to distinguish this number from zero.
  \end{enumerate}
\end{soln}
\end{document}